\documentclass[10pt,a4paper]{amsart}
\usepackage[margin=19mm]{geometry}
\usepackage{amsmath,amssymb,mathtools}
\usepackage[scr=rsfs,cal=euler]{mathalfa}
\usepackage{xfrac}
\usepackage[unicode,hidelinks,bookmarks=false]{hyperref}
\newcommand{\ie}{i.e.\ }
\newcommand{\cf}{cf.\ }
\newcommand{\resp}{respectively\ }
\newcommand{\Subsection}{\subsection}
\providecommand{\nobreakdash}{\nobreak\hbox{-}}
\setlength{\emergencystretch}{2em}
\multlinegap0pt
\usepackage{bm}
\usepackage{bbm}
\usepackage{dsfont}
\usepackage{xspace}
\usepackage{cancel}
\usepackage{mathtools}
\usepackage{cleveref}
\usepackage{amsthm}
\makeatletter
\@ifundefined{theorem}{\newtheorem{theorem}{Theorem}}{}
\makeatother
\title[On Structural Properties and Adjacency Spectrum of Coprime G]{On Structural Properties and Adjacency Spectrum of Coprime Graph of Integers}
\author{Subarsha Banerjee}
\date{Source: arXiv:2506.10583v2 (CC BY 4.0)}
\begin{document}
\maketitle

\noindent\emph{Source and licence.} On Structural Properties and Adjacency Spectrum of Coprime Graph of Integers. Version arXiv:2506.10583v2 (CC BY 4.0), distributed under the Creative Commons Attribution 4.0 International licence, as recorded in the arXiv licence metadata for this exact version and shown on the abstract page https://arxiv.org/abs/2506.10583v2. Full rights notice: licenses/arxiv-2506.10583-CC-BY-4.0.txt.\par\smallskip

\noindent\textit{Editorial scope.} Counted unit: the Theorem whose source label is \texttt{number} in the article (source0.tex lines 401--404, source label \texttt{number}), delivered together with the article's own complete original proof of it (source0.tex lines 406--510, one \texttt{proof} environment of 105 lines). No mathematics is written, repaired or re-derived: statement and proof are the article's own text line for line. Required context retained uncounted: none. Every definition, display and label of those lines is retained unchanged. Omitted from this package and not redistributed: 1--400, 405--405, 511--855. Nothing inside a retained excerpt is omitted or altered: every retained body is delivered exactly as the article's lines are written, so no omission receipt of any kind accompanies this package. Citations: the retained excerpts carry no citation command at all, so no bibliography entry of the article is transcribed and no reference is redistributed. Reference adaptation: none was needed and none was made; every reference inside the delivered unit resolves inside it, so no unresolved reference or citation is produced. Printed slips kept literal: no printed word, symbol, hypothesis or number of the counted statement or of its proof was altered. Layout and class: the article's own class is \texttt{article} with packages including inputenc, graphicx, amsmath, amssymb, amsthm, soul, tikz, float, hyperref, rotating, geometry, blkarray, cleveref, fancyhdr; the wrapper is the lane's pinned \texttt{amsart} editorial wrapper at 10pt on A4, which restates the theorem-like environments the retained text needs and adds the article's own macro definitions that the retained text needs (none), with extra wrapper packages (bm, bbm, dsfont, xspace, cancel, mathtools, cleveref). The delivered numbering is the unit's own. Assets: no retained span contains a figure environment, a graphics-inclusion command, a PDF-page inclusion, a table or a TikZ picture; no image file is copied into this package, the package declares no asset dependency and no image file is redistributed with it. Licence: version arXiv:2506.10583v2 of this article is distributed under the Creative Commons Attribution 4.0 International licence, as recorded in the arXiv licence metadata for this exact version and shown on the abstract page https://arxiv.org/abs/2506.10583v2. A full rights notice is delivered at licenses/arxiv-2506.10583-CC-BY-4.0.txt. Characters: every retained excerpt is pure ASCII, so the delivered engine, XeLaTeX with its default Latin Modern text font, reports no missing character.
	\begin{theorem}
		\label{number}
		The number of times $-1$ appears as an eigenvalue of the adjacency matrix of $TCG_n$ is at least equal to the cardinality of $\mathcal{D}$, where $\mathcal{D}$ is the set of prime numbers which are greater than $\frac{n}{2}$ but less than $n$. 
	\end{theorem}

	\begin{proof}
		Let $\mathcal{D}$ be the collection of all the prime numbers which are greater than $\frac{n}{2}$ but less than $n$. Let $p\in \mathcal{D}$. 
		Then $\gcd(p,a)=1$ for all $1\le a< p$ and $p<a\le n$.
		This implies that $p$ is adjacent to all the members of the graph $TCG_n$.
		Thus, the subgraph of $TCG_n$ induced by the set $\mathcal{D}\cup \{1\}$ is a complete graph on $|\mathcal{D}|+1$ vertices.
		We shall index the rows and columns of $A(TCG_n)$ as follows:
		We first index the vertex $1$, we then index all the vertices of $TCG_n$ which are present in $\mathcal{D}$, and finally we index all the remaining vertices of 
		$TCG_n$.
		Using the above indexing, we find that $A(TCG_n)$ takes the following form:
		
		\begin{equation}
			\label{Eq0}
			\begin{split}
				A=A(TCG_n)&=
				\\
				&
				\left[
				\begin{array}{cccccccccccccccccc}
					0 & 1 & 1& \ldots &\dots & \textbf{1}&1 & 1 & 1& \ldots&\ldots &\ldots& 1\\
					1 & 0 & 1& \ldots &\dots &\textbf{1}&1 & 1 & 1& \ldots &\ldots&\ldots& 1\\
					1& 1 & 0 &\ldots&\dots& \textbf{1}&1 & 1 & 1& \ldots &\ldots&\ldots& 1\\
					\ldots& \ldots &\ldots &\ldots&\ldots& \ldots &\ldots &\ldots&\ldots&\ldots&\ldots&\ldots&\ldots\\
					\ldots& \ldots &\ldots &\ldots&\ldots& \ldots &\ldots &\ldots&\ldots&\ldots&\ldots&\ldots&\ldots\\
					\ldots& \ldots &\ldots &\ldots&\ldots& \ldots &\ldots &\ldots&\ldots&\ldots&\ldots&\ldots&\ldots\\
					\textbf{1}& \textbf{1} &\textbf{1}&\ldots &\ldots& \textbf{0} &\textbf{1} & \textbf{1} & \textbf{1}& \ldots &\ldots&\ldots& \textbf{1}\\
					1 & 1 &1  & \ldots&\ldots&\textbf{1} & 0 & a_{12} &a_{13}& \ldots&\ldots&\ldots&a_{1m}\\
					1 & 1 &1 & \ldots&\ldots & \textbf{1}& a_{21} & 0 &a_{23}& \ldots&\ldots&\ldots&a_{2m}\\
					1& 1 & 1 &\ldots&\dots& \textbf{1}&a_{31} & a_{32} & 0& \ldots &\ldots&\ldots& a_{3m}\\
					\ldots& \ldots &\ldots &\ldots&\ldots& \ldots &\ldots &\ldots&\ldots& \ldots&\ldots&\ldots&\ldots \\
					\ldots& \ldots &\ldots &\ldots&\ldots& \ldots &\ldots &\ldots&\ldots& \ldots &\ldots&\ldots&\ldots\\
					\ldots& \ldots &\ldots &\ldots&\ldots& \ldots &\ldots &\ldots&\ldots& \ldots &\ldots&\ldots&\ldots\\
					\ldots& \ldots &\ldots &\ldots&\ldots& \ldots &\ldots &\ldots&\ldots& \ldots &\ldots&\ldots&\ldots\\
					\ldots& \ldots &\ldots &\ldots&\ldots& \ldots &\ldots &\ldots&\ldots&\ldots&\ldots&\ldots&\ldots\\
					1 & 1 & 1 &  \ldots& \dots&\textbf{1}& a_{m1} & a_{m2} &  a_{m3}&\ldots&\ldots&\ldots& 0\\
				\end{array}
				\right].
			\end{split}
		\end{equation}
		Here $m=n-1-|\mathcal{D}|.$
		
		
		Consider the matrix $A+I$, where $I$ is the identity matrix of order $n$.
		We have,
		
		\begin{equation}
			\label{Eq1}
			\begin{split}
				A+I&=
				\\
				&
				\left[
				\begin{array}{cccccccccccccccccc}
					1 & 1 & 1& \ldots &\dots & \textbf{1}&1 & 1 & 1& \ldots&\ldots &\ldots& 1\\
					1 & 1 & 1& \ldots &\dots &\textbf{1}&1 & 1 & 1& \ldots &\ldots&\ldots& 1\\
					1& 1 & 1 &\ldots&\dots& \textbf{1}&1 & 1 & 1& \ldots &\ldots&\ldots& 1\\
					\ldots& \ldots &\ldots &\ldots&\ldots& \ldots &\ldots &\ldots&\ldots&\ldots&\ldots&\ldots&\ldots\\
					\ldots& \ldots &\ldots &\ldots&\ldots& \ldots &\ldots &\ldots&\ldots&\ldots&\ldots&\ldots&\ldots\\
					\ldots& \ldots &\ldots &\ldots&\ldots& \ldots &\ldots &\ldots&\ldots&\ldots&\ldots&\ldots&\ldots\\
					\textbf{1}& \textbf{1} &\textbf{1}&\ldots &\ldots& \textbf{1} &\textbf{1} & \textbf{1} & \textbf{1}& \ldots &\ldots&\ldots& \textbf{1}\\
					1 & 1 &1  & \ldots&\ldots&\textbf{1} & 1 & a_{12} &a_{13}& \ldots&\ldots&\ldots&a_{1m}\\
					1 & 1 &1 & \ldots&\ldots & \textbf{1}& a_{21} & 1 &a_{23}& \ldots&\ldots&\ldots&a_{2m}\\
					1& 1 & 1 &\ldots&\dots& \textbf{1}&a_{31} & a_{32} & 1& \ldots &\ldots&\ldots& a_{3m}\\
					\ldots& \ldots &\ldots &\ldots&\ldots& \ldots &\ldots &\ldots&\ldots& \ldots&\ldots&\ldots&\ldots \\
					\ldots& \ldots &\ldots &\ldots&\ldots& \ldots &\ldots &\ldots&\ldots& \ldots &\ldots&\ldots&\ldots\\
					\ldots& \ldots &\ldots &\ldots&\ldots& \ldots &\ldots &\ldots&\ldots& \ldots &\ldots&\ldots&\ldots\\
					\ldots& \ldots &\ldots &\ldots&\ldots& \ldots &\ldots &\ldots&\ldots& \ldots &\ldots&\ldots&\ldots\\
					\ldots& \ldots &\ldots &\ldots&\ldots& \ldots &\ldots &\ldots&\ldots&\ldots&\ldots&\ldots&\ldots\\
					1 & 1 & 1 &  \ldots& \dots&\textbf{1}& a_{m1} & a_{m2} &  a_{m3}&\ldots&\ldots&\ldots& 1\\
				\end{array}
				\right].
			\end{split}
		\end{equation}
		
		We find that the first $\mathcal{D}+1$ rows of the matrix $A+I$ are identical.
		Hence we apply the row operation $R_i=R_i-R_1$ for $2\le i\le \mathcal{D}+1$.
		On applying the given row operations, the matrix $A+I$ reduces to the following matrix:
		
		\begin{equation}
			\label{Eq2}
			\left[
			\begin{array}{cccccccccccccccccccccccccccccccccccccccccccccc}
				1 & 1 & 1& \ldots &\dots & \textbf{1}&1 & 1 & 1& \ldots&\ldots &\ldots& 1\\
				0 & 0 & 0& \ldots &\dots &\textbf{0}&0 & 0 & 0& \ldots &\ldots&\ldots& 0\\
				0& 0 & 0 &\ldots&\dots& \textbf{0}&0 & 0 & 0& \ldots &\ldots&\ldots& 0\\
				\ldots& \ldots &\ldots &\ldots&\ldots& \ldots &\ldots &\ldots&\ldots&\ldots&\ldots&\ldots&\ldots\\
				\ldots& \ldots &\ldots &\ldots&\ldots& \ldots &\ldots &\ldots&\ldots&\ldots&\ldots&\ldots&\ldots\\
				\ldots& \ldots &\ldots &\ldots&\ldots& \ldots &\ldots &\ldots&\ldots&\ldots&\ldots&\ldots&\ldots\\
				\textbf{0}& \textbf{0} &\textbf{0}&\ldots &\ldots& \textbf{0} &\textbf{0} & \textbf{0} & \textbf{0}& \ldots &\ldots&\ldots& \textbf{0}\\
				\cline{1-13}
				&   &                      &   & \multicolumn{1}{c|}{}  &   &                        \\
				&   & \makebox[0pt]{\huge{$J^T$}}    &   & \multicolumn{1}{c|}{}  &&   &&\makebox[0pt]{\huge{$A(G')$}}    \\
				&   &                      &   & \multicolumn{1}{c|}{}  &   &
			\end{array} \;
			\right].
		\end{equation}
		Here $G'$ is the induced subgraph of $TCG_n$ on the set $\mathbb Z_n\setminus \{\mathcal{D}\cup \{1\}\}.$ 
		$A(G')$ denotes the adjacency matrix of $G'$, and $J$ denotes the all one matrix.
		We note that the matrix in \Cref{Eq2}  has $\mathcal{D}$ identical rows of the form: 
		$\begin{array}{ccccccccccccccccc}
			(0 & 0 & 0 &\dots & 0 &0 &0&0&0&\dots&0 &0 ).
		\end{array}$
		Thus rank of $A+I$ is less than $n-\mathcal{D}$. Using the Rank-Nullity Theorem, we find that  $0$ is an eigenvalue of $A+I$ with multiplicity at least $\mathcal{D}$. Hence, we conclude that $-1$ is an eigenvalue of $A$ with multiplicity at least $\mathcal{D}$.
		
		
	\end{proof}

\end{document}
